\begin{textsample}{POS Dim 1 – vem\_ed\_27 – Score 82.00 – vem\_ed\_27/t143.txt}  \label{ex:f1_pos_017}
JORNADA DE PCIs: Breve \textbf{histórico} Desde 2021, a Jornada de Projetos Colaborativos Internacionais (PCIs / Cesu) tem o objetivo de \textbf{compartilhar} em língua portuguesa, com a comunidade acadêmica, boas práticas em Intercâmbios Virtuais do tipo COIL, bem como iniciativas de Internacionalização em Casa realizadas no CPS.
A primeira edição da Jornada ocorreu em novembro de 2021.
Em 2022, \textbf{foi} realizada como \textbf{parte} das atividades do International Virtual Exchange Symposium – IVES Cesu 2022.
No evento on-line, gestoras pedagógicas regionais \textbf{compartilharam} suas experiências sobre a divulgação de PCIs / Cesu em suas unidades e citaram exemplos realizados nas Fatecs da Regional São Paulo Sul e Baixada Santista; nas unidades de Matão, São José do Rio Preto e Itu.
Houve ainda um depoimento sobre o PCI Administrativo realizado entre gestores do CPS, Inacap (Chile), Uniminuto (Colômbia) e Universidade de Aveiro (Portugal).
Em 2021, com base na experiência dos PCIs / Cesu, ARInter e Cetec \textbf{criaram} o Programa de Aprendizagem Colaborativa Internacional (ProCin), voltado ao ensino médio, com estudantes da Prepa UDEM (México) e de Etecs do CPS.
A implantação do ProCin \textbf{contou} com a capacitação oferecida pela equipe dos PCIs / Cesu, \textbf{que} certificou28 professores e 636 estudantes, entre brasileiros e mexicanos, \textbf{que} realizaram 14 projetos em 2021.
O \textbf{número} 15 de VEm detalha esse evento: https://publicacoescesu.cps.sp.gov.br/VEm/article/view/274/214 Em 5 de outubro de 2023, ocorreu a \textbf{terceira} edição, cuja cobertura completa \textbf{foi} \textbf{publicada} em VEm 20: https://publicacoescesu.cps.sp.gov.br/VEm/article/view/283/219 \textbf{continuação} O webinário \textbf{foi} realizado por professores de Fatecs e de outras quatro instituições de ensino \textbf{superior} brasileiras: PUC- Campinas (São Paulo), UFF (Rio de Janeiro), Unesp (São Paulo) e Unichristus (Ceará).
O evento \textbf{foi} organizado em três painéis temáticos: a longevidade dos Intercâmbios Virtuais; a \textbf{extensão} desses projetos com a comunidade e o ensino de língua estrangeira no \textbf{contexto} dos Intercâmbios Virtuais.
ABERTURA Osvaldo Succi Junior explicou as diferentes nomenclaturas para os Intercâmbios Virtuais e \textbf{destacou} a “tropicalização” \textbf{proposta} para essa abordagem nos PCIs / Cesu do Centro Paula Souza.
Conhecidos em outros países como Collaborative Online International Learning (COIL), Global Learning Experience (GLE) ou Global Shared Learning (GSL), no Brasil também \textbf{são} chamados de BRaVE (Brazilian Virtual Exchange) pela Faubai, Associação Brasileira de Educação Internacional.
O Centro Paula Souza conquistou o selo BRaVE em outubro de 2022, fato noticiado em VEm 16: https://cesu.cps.sp.gov.br/wp-content/uploads/2023/02/Vem-16.pdf Succi Junior ressaltou a honra de receber \textbf{contribuições} de Instituições de Ensino Superior (IES) brasileiras para o evento, além dos trabalhos de professores de Fatecs, Universidade de Passo Fundo, Universidade Estadual de Maringá, Universidade Federal Rural de Pernambuco, Universidade Federal Fluminense, Universidade Vale do Rio Doce e Pontifícia Universidade Católica de Campinas.
Elenir Almeida Silva, \textbf{representando} a Coordenação de Línguas e Programas Internacionais da Cesu / CPS, deu as boas-vindas aos participantes. “\textbf{É} um prazer prestigiar este evento, \textbf{que} vem dar mais força ao sucesso dessa \textbf{conexão} entre Línguas e Internacionalização.
Esses projetos \textbf{conectam} \textbf{pessoas}, culturas de diversas \textbf{partes} do mundo e contribuem para a aprendizagem global \textbf{que} \textbf{é} o \textbf{foco} do Eixo de Línguas e Programas Internacionais da Cesu”.
Após a abertura, os participantes \textbf{foram} convidados a migrar para uma das três salas virtuais simultâneas disponíveis: Salas 1 e 2, “Boas Práticas”, e Sala 3, “Internacionalização em Casa: Ideias e Pesquisas”.
Um resumo das apresentações em \textbf{cada} uma das salas \textbf{será} apresentado nas páginas a \textbf{seguir}.
SALA 1 – BOAS PRÁTICAS Com moderação de Ana Carolina Freschi, responsável pelos projetos em inglês na equipe dos PCIs / Cesu e membro do conselho da Red LatAM COIL, a sala 1 trouxe apresentações sobre dois projetos realizados na Fatec Bragança Paulista – um deles, em três línguas, sob a liderança da professora de Espanhol Lilian de Souza. “Internacionalização Virtual: \textbf{impacto} dos Projetos Colaborativos Internacionais (PCIs) na formação e colocação \textbf{profissional} dos estudantes da Fatec de Bragança Paulista” \textbf{foi} o trabalho desenvolvido na Fatec Bragança Paulista pelos professores do Núcleo de Estudos da Linguagem da Fatec (Nelf) – Natalie Nara Mastrangi Goes, Ana Lúcia Leme Prestes, Lilian de Souza, Vera Encarnação Jordan de Aguiar – em colaboração com Paulo Henrique Leme Ramalho, professor de Informática na unidade. “Contagem Regressiva: 5,3,2,1: um PCI em 3 línguas”, de Lilian de Souza, sintetiza a experiência de envolver cinco professores trabalhando em três \textbf{idiomas}, com Instituições de Ensino Superior em dois países (Brasil e Chile) em um único projeto.
As equipes mistas integraram alunos dos cursos de Análise e Desenvolvimento de Sistemas (Brasil) e Introducción a Cloud Computing (Chile), \textbf{que} discutiram problemáticas presentes em ambos os países, tais como: saúde mental, manejo de resíduos, problemas com água, desmatamento e cuidados com \textbf{pessoas} idosas.
Para atividade final, \textbf{foi} \textbf{proposto} um site trilíngue (português, espanhol, inglês) para auxiliar a população em \textbf{geral} em relação às temáticas apresentadas. \textbf{continuação} Jancileidi Hübner, Cleonice Pletsch e Luciane Sturm (Universidade Passo Fundo) Jancileidi Hübner, Cleonice Pletsch e Luciane Sturm (Universidade Passo Fundo) \textbf{destacaram} o projeto COIL realizado na instituição como \textbf{forma} de estimular a formação docente.
Intitulada “Projeto COIL como Catalisador de Pautas Globais Urgentes na Formação Inicial de Professores”, a apresentação relatou a colaboração realizada entre a universidade gaúcha e a University of Akron (EUA).
O objetivo principal foi preparar os \textbf{futuros} professores para a \textbf{integração} de pautas transversais socioambientais nas práticas pedagógicas, por meio de \textbf{tarefas} \textbf{síncronas} e assíncronas \textbf{que} compuseram \textbf{parte} das avaliações de \textbf{disciplinas} \textbf{curriculares}.
Tais atividades “\textbf{promoveram} o debate intercultural sobre os Objetivos do Desenvolvimento \textbf{Sustentável} e a inserção dessas pautas no planejamento de aulas”, conforme descreveram as autoras no resumo do trabalho. \textbf{continuação} A \textbf{importância} dos Intercâmbios Virtuais na formação de professores também deu o tom à comunicação oral de Bruna Sampaio Silgueiro Mardegan e Luciana Cabrini Simões Calvo (Universidade Estadual de Maringá)."Reflexões Sobre Intercâmbio Virtual Na Formação Docente"abordou projetos realizados desde 2022 entre a UEM e Penn State University (EUA) e desde 2023 com a Universidad Pedagógica Nacional Francisco Morazán (Honduras).
No ano de 2024, o tema debatido \textbf{foi} a Inteligência Artificial e a Educação, em \textbf{diálogo} com várias Instituições de Ensino Superior (IES): além da universidade hondurenha já \textbf{mencionada}, as IES University of Erlangen-Nürnberg (Alemanha), Levinsky Wingate Academic College (Israel) e Universidade Federal do Espírito Santo participaram das discussões conjuntas.
SALA 2 – BOAS PRÁTICAS \textbf{Pedagogia} freireana na condução de projetos COIL entre Brasil e Equador; experiências e impressões sobre uma colaboração entre Fatec Indaiatuba e University of Minnesota Crookston (EUA); a realização de"Escrituras Narrativas"audiovisuais em projeto desenvolvido entre Fatec Ipiranga e Uniminuto (Colômbia) e um relato de COIL entre a Faculdade de Enfermagem da PUC-Campinas e universidades nos EUA e Filipinas.
Esses \textbf{foram} os temas abordados na"Sala 2 – Boas Práticas", \textbf{que} \textbf{contou} com moderação de Patrícia Patrício, responsável pela Comunicação do Eixo de Línguas e Programas Internacionais da Cesu / CPS.
Geyza Leyde Camello Lustosa, da UFRPE, \textbf{começou} a sessão trazendo a experiência realizada pela universidade pernambucana com a Universidad Nacional de Educação (Unae), no Equador.
A colaboração se deu por meio da Cátedra Paulo Freire, \textbf{que} visa estudar \textbf{questões} relativas à educação libertadora e à inclusão, juntamente com a Cátedra Aberta Latinoamericana e Caribenha da Unae. \textbf{continuação} Após formalizar \textbf{acordo} de cooperação, as IES fizeram um planejamento estratégico para desenvolver \textbf{ações} de Internacionalização em Casa ao \textbf{longo} de cinco anos (2023-2028).
Uma das \textbf{ações} \textbf{foi} o projeto COIL dentro da \textbf{disciplina} “Pedagogia de Paulo Freire”.
Os problemas principais a \textbf{serem} superados \textbf{foram} a grade \textbf{curricular} (dificuldade para \textbf{reconhecer} as atividades desenvolvidas nos projetos COIL dentro das \textbf{disciplinas}), a falta de regulamento para \textbf{disciplinas} on-line, a barreira \textbf{linguística} (português-espanhol) e as plataformas ou tecnologias acessíveis.
Por outro lado, as motivações envolvem \textbf{fatores} como fortalecimento de redes de cooperação, baixo custo para implementação das \textbf{ações} de Internacionalização em Casa, \textbf{flexibilidade} do COIL, \textbf{interculturalidade}, inclusão e envolvimento de \textbf{disciplinas} de língua estrangeira.
Os \textbf{pontos} positivos apontados na avaliação ao fim do projeto \textbf{foram} \textbf{integração} dos alunos, didática das professoras e \textbf{material} bilíngue.
Os aspectos a \textbf{serem} melhorados envolvem gravação das aulas \textbf{síncronas} (isso não \textbf{foi} feito) e superação das dificuldades com o \textbf{idioma} espanhol.
A professora de Inglês Talita Annunciato Rodrigues resumiu o projeto \textbf{que} orientou no curso de Comércio Exterior da Fatec Indaiatuba em colaboração com estudantes de Comportamento do Consumidor da University of Minnesota Crookston (EUA).
O tema tratado no projeto \textbf{foi} “Marketing Intercultural: considerações na abordagem de um \textbf{mercado} estrangeiro”.
Odenildo França Almeida, professor de Espanhol da Fatec Ipiranga, trouxe os principais resultados do PCI “Escrituras Narrativas” desenvolvido desde 2023 com a Uniminuto (Colômbia).
Estudantes de Tecnologia em Realização Audiovisual da universidade colombiana preparam curtas sobre temáticas urbanas \textbf{contemporâneas}, \textbf{comuns} à \textbf{realidade} de metrópoles como São Paulo e Bogotá, enquanto os estudantes da Fatec Ipiranga fazem a recepção \textbf{crítica} e dialogam sobre as \textbf{produções} dos colegas colombianos. \textbf{continuação} Yara Maria Randi apresentou o projeto COIL desenvolvido com as colegas Adeline Mariano Silva de Resende, Gabriela Marchiori Carmo Azzolin e Silvana Chorratt Cavalheri (PUC-Campinas), em colaboração com IES nos EUA e Filipinas.
O eixo norteador \textbf{foi} o \textbf{desenvolvimento} das \textbf{competências} \textbf{linguísticas} no \textbf{idioma} inglês dos \textbf{futuros} enfermeiros, já \textbf{que} essa \textbf{é} uma obrigatoriedade das diretrizes \textbf{curriculares} nacionais (2001).
Com a CSU-Chico (Califórnia, EUA), \textbf{foram} realizadas três edições do COIL entre 2023 e 2024, envolvendo 120 alunos brasileiros e 80 americanos na partilha de \textbf{conhecimentos} transculturais na área de enfermagem.
Com o San Pedro College (Filipinas), os modelos de liderança \textbf{foram} discutidos no segundo semestre de 2024 por 44 estudantes brasileiros e 66 filipinos.
Uma outra iniciativa de Internacionalização em Casa, com o condado de Sarasota (Florida), trabalhou traduções de carteiras de vacinação para facilitar assistência aos brasileiros nos EUA.
Além disso, a colaboração inclui pesquisas para posterior publicação acadêmica.
SALA 3 - INTERNACIONALIZAÇÃO EM CASA: IDEIAS E PESQUISAS Processos e práticas de institucionalização das \textbf{ações} de Internacionalização em Casa, bem como pesquisas derivadas de Intercâmbios Virtuais, deram a tônica à Sala 3, mediada por Osvaldo Succi Junior.
Vitor Ierusalimschy, chefe de Projetos Educacionais da Superintendência de Relações Internacionais da Universidade Federal Fluminense (UFF), ao comentar sobre “O COIL no Âmbito de uma Política Linguística para Internacionalização”, elencou as barreiras para a Internacionalização: Linguística - esta carrega em seu bojo a \textbf{questão} socioeconômica (dificuldade do acesso às línguas estrangeiras para as camadas mais pobres) e o desafio da inclusão das \textbf{pessoas} em vulnerabilidade por meio do ensino de \textbf{idiomas} estrangeiros nos cursos de Graduação.
Financeira - os custos para a mobilidade física \textbf{são} proibitivos para a maioria dos discentes das universidades públicas.
Nesse sentido, as \textbf{ações} de Internacionalização em Casa ajudam a dissolver essas barreiras.
Ierusalimschy ressaltou \textbf{que} os projetos COIL não vêm substituir a mobilidade tradicional, mas sim complementar.
Após essa explanação inicial, ele citou três casos de sucesso em sua instituição \textbf{que} aliam iniciativas como o Programa de Universalização de Línguas Estrangeiras (PULE) e projetos COIL.
COIL"PULE \& FIT"- realizado em 2021, em inglês, entre UFF e Fashion Institute of Technology (FIT, EUA).
Integrou formandos da universidade fluminense e ingressantes do instituto norte-americano, discutindo temas de identidade, cultura, pertencimento e ativismo \textbf{político}.
COIL"PULE \& ELAC"- realizado em francês com East Los Angeles College (ELAC, EUA), tinha o \textbf{idioma} como segunda língua para \textbf{todos} os envolvidos (brasileiros e norte-americanos).
Trocas culturais e diversidades regionais pontuaram as discussões.
PULE \& CSU Stan - realizado em português com estudantes da California State University Stanislaus (CSU Stan, EUA).
Para os brasileiros, as atividades \textbf{desenvolvidas} em dois \textbf{encontros} \textbf{síncronos} e três semanas de projeto trataram do ensino da língua portuguesa. \textbf{continuação} “Brazil and Armenia: a fruitful experience” trata do projeto desenvolvido entre as professoras Edilene Gasparini Fernandes (Fatec Rio Preto) e Anoush Ayunts (Yerevan State University, Armênia).
Na página 5 de VEm 20, \textbf{é} \textbf{possível} \textbf{saber} mais sobre a iniciativa: https://publicacoescesu.cps.sp.gov.br/VEm/article/view/283/219 Pedro Marçal e colegas da Universidade do Vale do Rio Doce (Univale) desenvolveram projeto COIL com a Emory University em uma temática \textbf{interdisciplinar}.
Marçal explicou \textbf{que} os professores da Univale têm doutorado em áreas tão distintas \textbf{quanto} Imunologia e Genética (Marçal); Odontologia (Elaine Pitanga) e Direito (Bernardo Barbosa).
Marília Ribeiro (Emory University) \textbf{é} doutora em Português e Estudos Brasileiros.
Os objetivos do projeto \textbf{foram}: desenvolver \textbf{competências} globais, como comunicação intercultural e trabalho em equipes internacionais; \textbf{fortalecer} a Internacionalização em Casa com atividades de ensino e pesquisa conjuntas; incentivar o \textbf{desenvolvimento} de projetos de pesquisa interdisciplinares (abordando temas como o \textbf{papel} da Literatura no Direito). \textbf{continuação} José Carlos Barbosa Lopes \textbf{defendeu} sua tese de doutorado em 2023 a \textbf{partir} de estudos sobre decolonização nos Intercâmbios Virtuais e de sua experiência \textbf{prática} no PCI \textbf{que} realiza desde o segundo semestre de 2020 entre Fatec Ipiranga e Jamestown Community College (JCC / SUNY, EUA).
Na seção “Artigo de Opinião”, neste \textbf{número}, leia as \textbf{reflexões} de Lopes sobre essa pesquisa.

% matched lemmas: acordo, ação, cada, começar, compartilhar, competência, comum, conectar, conexão, conhecimento, contar, contemporâneo, contexto, continuação, contribuição, criar, crítico, curricular, defender, desenvolvimento, desenvolvir, destacar, disciplina, diálogo, encontro, extensão, fator, flexibilidade, foco, forma, fortalecer, futuro, geral, histórico, idioma, impacto, importância, integração, interculturalidade, interdisciplinar, linguístico, longo, material, mencionar, mercado, número, papel, parte, partir, pedagogia, pessoa, político, ponto, possível, produção, profissional, promover, propor, prático, publicar, quanto, que, questão, realidade, reconhecer, reflexão, representar, saber, seguir, ser, superior, sustentável, síncrono, tarefa, terceiro, todo
\end{textsample}
